\section{Poisson summation and logarithmic second moments}
\label{analytic:section}
\label{sec:analytic}

We retain the finite family $\mathcal F$, the fixed set $S$, and the
unitary normalization from Lemmas~\ref{setup:reduction} and
\ref{setup:local-factors}. In particular, $|\lambda_F(n)|\leq\tau(n)$,
all local parameters have modulus at most one, and reduction of a
squarefree twisting parameter to an admissible one does not increase its
absolute value. All constants in this section are uniform in
$F\in\mathcal F$.

Our first aim is to bound the second moments of the twist functions and
their localized Dirichlet polynomials with only logarithmic losses.
For a compactly supported smooth function $G$ on $(0,\infty)$, write
\[
 P_d(N,t;G)=\sum_{n\geq1}
 \lambda_F(n)\chi_d(n)n^{-1/2-it}G(n/N).
\]
All dyadic parameters below belong to $\{1,2,4,\ldots\}$.

\begin{proposition}[Logarithmic second moments]\label{analytic:moment}
There are fixed positive integers $A,B$ such that, for every dyadic
$M\geq1$, every $F\in\mathcal F$, every $t\in\mathbb R$, and
$1/2\leq\sigma\leq2$,
\begin{equation}\label{analytic:L-moment}
 \frac1M\sum_{\substack{d\ \mathrm{admissible}\\M\leq|d|<2M}}
 |L_{F^{(d)}}(\sigma+it)|^2
 \ll_{\mathcal F}\log^B(3M)(1+|t|)^A.
\end{equation}
For arbitrary signed squarefree twisting parameters and $H\geq1$,
\begin{equation}\label{analytic:prefix-moment}
 \sum_{\substack{0<|h|\leq H\\h\ \mathrm{squarefree}}}
 |L_{F^{(h)}}(\sigma+it)|^2
 \ll_{\mathcal F}H\log^B(3H)(1+|t|)^A.
\end{equation}
\end{proposition}

\begin{proposition}[Dirichlet-polynomial second moments]
\label{analytic:polynomial-moment}
With the same fixed integers $A,B$, every fixed smooth $G$ of compact
support in $(0,\infty)$ satisfies
\begin{equation}\label{analytic:polynomial-bound}
 \frac1M\sum_{\substack{d\ \mathrm{admissible}\\M\leq|d|<2M}}
 |P_d(N,t;G)|^2
 \ll_{\mathcal F,G}
 \log^B(3M)(1+|t|)^A(1+N/M)^A
 \qquad(N>0).
\end{equation}
The bound is uniform for fixed uniformly smooth families of weights
with common compact support.
\end{proposition}

We prove these estimates by induction on the twisting height $M$.
Poisson summation turns a polynomial mean square into a sum of products
of twist $L$-functions at dual squarefree parameters. To put those
parameters below the current height, we first enlarge the family by
replacing $d$ with $dp^2$, where $p\asymp P$. The new averaging height is
$U=MP^2$. A weight with compact Fourier support then restricts the dual
parameters to size $O(N^2/(MP^2))$, which is at most $M/2$ for a
sufficiently large $P$. Recovering the original squarefree average costs
$O(P\log(2P))$. The nonzero-frequency estimate gains $P^{-2}$, so after
recovery its recursive coefficient contains $N\log(2P)/(MP)$.
We choose $P$ as a sufficiently large fixed power of
$\log(3M)(1+N/M)(1+|t|)$, with a large fixed leading constant, to make
this coefficient small. The approximate functional equation and a strip
estimate then close the induction for the $L$-functions themselves.
The compact Fourier support construction and prime-square inflation
are related to Li's treatment of quadratic twists \cite{Li2022}.

The arithmetic preparation below keeps its arithmetic exponents independent of the
number of weight derivatives. Once the induction is closed, we sum all
dual frequencies to compare short polynomial averages with the integer
model. The polynomial moment bound will also control the unmollified
terminal series in Section~\ref{sec:weights}.

\subsection{The normalized Gauss sums}

Put $e(x)=\exp(2\pi i x)$. For an odd positive integer $l$, set
\[
 \tau_h(l)=\sum_{a\bmod l}\left(\frac a l\right)e(ah/l),\qquad
 \epsilon_l=\begin{cases}1&l\equiv1\pmod4,\\i&l\equiv3\pmod4,\end{cases}
 \qquad B_h(l)=\frac{\tau_h(l)}{\epsilon_l\sqrt l}.
\]
We use $B_h(1)=1$. The square root is positive.
These are the normalized quadratic Gauss sums of
\cite[Section~2.2, Lemma~2.3]{Soundararajan2000}: in that notation,
$B_h(l)=G_h(l)/\sqrt l$. We give the formulas and their proof in the
present normalization.

\begin{lemma}[Exact Gauss-sum formulas]\label{analytic:gauss}
For fixed $h$, the function $l\mapsto B_h(l)$ is multiplicative on odd
positive integers, and
\begin{equation}\label{analytic:phase}
 B_{-h}(l)=\left(\frac{-1}{l}\right)B_h(l),\qquad
 \epsilon_l B_h(l)=\frac{1+i}{2}B_h(l)+\frac{1-i}{2}B_{-h}(l).
\end{equation}
If $h\ne0$, $p$ is odd, $t=v_p(h)$, and $v\geq1$, then
\begin{equation}\label{analytic:prime-powers}
 B_h(p^v)=
 \begin{cases}
 p^{v/2}(1-p^{-1}),&v\text{ even},\ v\leq t,\\
 -p^{v/2-1},&v\text{ even},\ v=t+1,\\
 p^{(v-1)/2}\left(\dfrac{h/p^{v-1}}p\right),
       &v\text{ odd},\ v=t+1,\\
 0,&\text{otherwise}.
 \end{cases}
\end{equation}
Consequently $B_h(p^v)=0$ for $v>t+1$ and
$|B_h(p^v)|\leq p^{v/2}$. If $t=0$, only $v=1$ can survive and
$B_h(p)=(h/p)$. If $t=1$, only $v=2$ can survive and
$B_h(p^2)=-1$.
\end{lemma}

\begin{proof}
For coprime odd $l_1,l_2$, the Chinese remainder theorem, followed by a
change of variable in each Gauss sum, gives
\[
 \tau_h(l_1l_2)=
 \left(\frac{l_1}{l_2}\right)\left(\frac{l_2}{l_1}\right)
 \tau_h(l_1)\tau_h(l_2).
\]
Quadratic reciprocity says that the displayed product of symbols is
$\epsilon_{l_1l_2}/(\epsilon_{l_1}\epsilon_{l_2})$. This proves
multiplicativity. Substituting $a\mapsto-a$ proves the first identity
in \eqref{analytic:phase}; the second follows by considering the two
residue classes of $l$ modulo four.

For even $v$, the Jacobi symbol modulo $p^v$ is the indicator of the
units. Thus
\[
 \tau_h(p^v)=\sum_{a\bmod p^v}e(ah/p^v)
             -\sum_{b\bmod p^{v-1}}e(bh/p^{v-1}),
\]
which equals $p^v-p^{v-1}$ when $p^v\mid h$, equals $-p^{v-1}$ when
$p^{v-1}\Vert h$, and otherwise vanishes. Here $\epsilon_{p^v}=1$.
For odd $v$, write $a=y+pz$ with $y$ modulo $p$ and $z$ modulo
$p^{v-1}$. The $z$-sum vanishes unless $p^{v-1}\mid h$; in that case
\[
 \tau_h(p^v)=p^{v-1}\sum_{y\bmod p}
       \left(\frac yp\right)e\big(y(h/p^{v-1})/p\big).
\]
This is zero if $p^v\mid h$. Otherwise the quadratic Gauss-sum formula
makes it
$p^{v-1}\epsilon_p\sqrt p\,\bigl(\frac{h/p^{v-1}}p\bigr)$.
Since $\epsilon_{p^v}=\epsilon_p$ for odd $v$, division by
$\epsilon_{p^v}p^{v/2}$ gives \eqref{analytic:prime-powers}.
\end{proof}

For later use, fix once and for all the arithmetic exponent
\begin{equation}\label{analytic:arithmetic-exponent}
 C_{\rm ar}=256.
\end{equation}
This exponent will not change when we require more derivatives of a
weight or more decay in an imaginary direction.

\begin{lemma}[Two-variable Euler factorization]\label{analytic:euler}
Let $r$ be positive and odd, and write $h=h_1h_2^2\ne0$, where $h_1$ is
signed squarefree and $h_2\geq1$. For $\nu\in\{1,-1\}$, the series
\[
 \mathcal G_{h,r}^{\nu}(s_1,s_2)=
 \sum_{\substack{n_1,n_2\geq1\\n_1,n_2\ \mathrm{odd}}}
 \frac{\lambda_F(n_1)\lambda_F(n_2)B_{\nu h}(rn_1n_2)}
      {n_1^{s_1}n_2^{s_2}}
\]
initially converges in a right half-plane and continues to
$\Re s_1,\Re s_2>1/2$ as
\begin{equation}\label{analytic:euler-factorization}
 \mathcal G_{h,r}^{\nu}(s_1,s_2)=
 L_{F^{(\nu h_1)}}(s_1)L_{F^{(\nu h_1)}}(s_2)
 \mathcal Q_{h,r}^{\nu}(s_1,s_2).
\end{equation}
For $0<b\leq1/4$ and $\Re s_i\geq1/2+b$,
\begin{equation}\label{analytic:euler-bound}
 |\mathcal Q_{h,r}^{\nu}(s_1,s_2)|
 \ll_{\mathcal F} b^{-C_{\rm ar}}r^{C_{\rm ar}}
                       \tau(h_2)^{C_{\rm ar}}.
\end{equation}
The function $\mathcal Q$ is holomorphic in these open half-planes.
\end{lemma}

\begin{proof}
For an odd prime put $e_p=v_p(r)$ and $t_p=v_p(h)$. Multiplicativity
expresses the series as the product of the local sums
\[
 T_p(s_1,s_2)=\sum_{a,b'\geq0}
 \lambda_F(p^a)\lambda_F(p^{b'})
 B_{\nu h}(p^{e_p+a+b'})p^{-as_1-b's_2}.
\]
The symbol $b'$ here is an integer exponent, distinct from the real
shift $b$. No division by $B_{\nu h}(p^{e_p})$ is made; that quantity
may vanish. Write $D_p^*(s)$ for the inverse local factor of the
elliptic curve $F^{(\nu h_1)}$. The correction factor at $p$ is
$T_p(s_1,s_2)D_p^*(s_1)D_p^*(s_2)$. At two it is just
$D_2^*(s_1)D_2^*(s_2)$, because the summation indices are odd.

There are three types of odd prime. If $p\notin S$ and $p\nmid rh$,
put $\chi=(\nu h/p)=(\nu h_1/p)$ and $x_i=p^{-s_i}$. The local
formulas give exactly
\[
 T_p=1+\lambda_F(p)\chi(x_1+x_2),\qquad
 D_p^*(s_i)=1-\lambda_F(p)\chi x_i+x_i^2.
\]
Multiplication cancels both terms of total degree one. The sum of the
absolute values of all polynomial coefficients before cancellation is
at most $5\cdot4\cdot4=80$. Therefore the correction is
$1+O(80p^{-1-2b})$, uniformly in the imaginary parts.

If $p\notin S$, $p\nmid r$, and $t_p=1$, the twisting curve has a
trivial local factor at $p$. By \eqref{analytic:prime-powers},
\[
 T_p=1-\lambda_F(p^2)(x_1^2+x_2^2)-\lambda_F(p)^2x_1x_2
     =1+O(10p^{-1-2b}).
\]
In particular these primes do not produce a cost for every prime
factor of the squarefree part $h_1$.

The remaining odd primes lie in $S$ or divide $rh_2$. The support in
\eqref{analytic:prime-powers} gives $e_p+a+b'\leq t_p+1$ for every
nonzero summand, except that $B_{\nu h}(1)=1$ is already covered when
$e_p=a=b'=0$. Hence
\[
 |T_p|\leq p^{e_p/2}
       \sum_{a+b'\leq t_p+1}(a+1)(b'+1)
 \leq p^{e_p/2}(t_p+3)^4.
\]
Each inverse local factor has modulus at most
$(1+p^{-1/2})^2$, so the product of the two costs at most $9$.
Writing $u_p=v_p(h_2)$, we have $t_p\leq2u_p+1$.
If $u_p\geq1$, the bound
$9(2u_p+4)^4\leq(u_p+1)^{17}$ absorbs this cost into a fixed divisor
power, with $p^{e_p/2}$ absorbed by $r$.
If $u_p=0$ and $e_p\geq1$, the cost is at most
$9\cdot4^4p^{e_p/2}\leq p^{8e_p}$ for $p\geq3$.
If $u_p=e_p=0$, this is one of the fixed primes of $S$, whose cost
is a fixed constant. The factor at two is bounded by
$(1+2^{-1/2})^4$.

The nonexceptional product is bounded by
\[
 \exp\left(80\sum_p p^{-1-2b}\right)
 \leq\zeta(1+2b)^{80}\ll b^{-80}.
\]
Combining the estimates proves \eqref{analytic:euler-bound} with the
stated $C_{\rm ar}$. The nonexceptional correction product converges
normally on compact subsets of $\Re s_i>1/2$; all exceptional factors
are polynomials. Multiplication by the entire twist $L$-functions
therefore gives the claimed continuation. This argument never divides
by a global $L$-value.
\end{proof}

\subsection{A weighted Poisson estimate}

The Euler factorization identifies the twist functions that occur at
nonzero frequencies. We now turn it into a weighted comparison formula:
the zero frequency gives the integer residue average, and the remaining
frequencies are bounded by those twist functions. This is a quadratic
Poisson method used in \cite[Section~2.4, Lemma~2.6]{Soundararajan2000},
with the weights and normalization specified below.

Fix a compact interval $[a_0,a_1]\subset(0,\infty)$ for the first two
scaled variables. The third variable ranges over $\mathbb R$. For an
integer $D\geq0$, a convenient finite seminorm is
\[
 \|W\|_{D,*}=\max_{\alpha+\beta+\gamma\leq D}
 \sup_{y_1,y_2,x}(1+|x|)^D
 \left|\partial_{y_1}^{\alpha}\partial_{y_2}^{\beta}
                    \partial_x^\gamma W(y_1,y_2,x)\right|.
\]
The functions under consideration vanish outside $[a_0,a_1]^2$ in the
first two variables and have the indicated finite decay and smoothness
in the third. Requiring a larger fixed $D$ below causes no change to
$C_{\rm ar}$.

\begin{lemma}[Poisson estimate]\label{analytic:poisson}
Let $U>0$, let $r$ be odd and positive, and let $N_i$ be bounded below
by a fixed positive constant. Fix $0<b\leq1/4$. In the expression
\begin{equation}\label{analytic:poisson-average}
 \frac1U\sum_{m\in\mathbb Z}
 \sum_{\substack{n_1,n_2\geq1\\n_1,n_2\ \mathrm{odd}}}
 \frac{\lambda_F(n_1)\lambda_F(n_2)}
 {(n_1n_2)^{1/2}n_1^{it_1}n_2^{it_2}}
 \left(\frac m{rn_1n_2}\right)
 W(n_1/N_1,n_2/N_2,m/U),
\end{equation}
Poisson summation gives the zero frequency, namely the complete-residue
average integrated in $x$, together with the nonzero frequencies.
For every fixed $J_1>0$, a sufficiently large fixed seminorm order $D$
bounds the absolute contribution of frequency $h=h_1h_2^2\ne0$ by
\begin{align}
 &C_{J_1,\mathcal F}\|W\|_{D,*}
 (N_1N_2)^{-1/2+b}b^{-C_{\rm ar}}r^{C_{\rm ar}}
 \tau(h_2)^{C_{\rm ar}}
 \left(1+\frac{|h|U}{rN_1N_2}\right)^{-J_1}
 \nonumber\\[-2pt]
 &\hspace{8mm}\cdot
 \sum_{\nu=\pm1}\int_{\mathbb R^2}
 \frac{\prod_{i=1}^2
 |L_{F^{(\nu h_1)}}(1/2+b+i(t_i+v_i))|}
 {(1+|v_1|+|v_2|)^{J_1}}\,dv_1\,dv_2.
 \label{analytic:nonzero-mode}
\end{align}
If the Fourier transform of $W$ in $x$ is supported in a fixed interval
$[-C_W,C_W]$, all modes with
$|h|>C_Wa_1^2rN_1N_2/U$ vanish exactly.
\end{lemma}

\begin{proof}
Our Fourier convention is
$\widehat W(y_1,y_2,\xi)=\int_{\mathbb R}W(y_1,y_2,x)e(-x\xi)\,dx$.
Applying Poisson summation in each residue class modulo $l=rn_1n_2$
gives exactly
\[
 \frac1U\sum_m\left(\frac ml\right)W(y_1,y_2,m/U)
 =\sum_{h\in\mathbb Z}\frac{\tau_h(l)}l
                         \widehat W(y_1,y_2,hU/l).
\]
Thus the normalized frequency factor is
\[
 \frac{1}{\sqrt{n_1n_2}}\frac{\tau_h(rn_1n_2)}{rn_1n_2}
 =\frac{\epsilon_{rn_1n_2}B_h(rn_1n_2)}
       {\sqrt r\,n_1n_2}.
\]
Use \eqref{analytic:phase} to replace the numerator by a fixed linear
combination of $B_h$ and $B_{-h}$.

Put $\rho=hU/(rN_1N_2)$ and
\[
 V_h(y_1,y_2)=\widehat W(y_1,y_2,\rho/(y_1y_2)),\qquad
 \widetilde V_h(z_1,z_2)=
 \int_0^\infty\!\int_0^\infty
 V_h(y_1,y_2)y_1^{z_1}y_2^{z_2}
                    \frac{dy_1dy_2}{y_1y_2}.
\]
On any fixed bounded range of real parts, Fourier integration by parts
in $x$, followed by Mellin integration by parts in $y_1,y_2$, yields
\begin{equation}\label{analytic:mellin-kernel}
 |\widetilde V_h(z_1,z_2)|
 \ll_{J_1}\|W\|_{D,*}(1+|\rho|)^{-J_1}
                 (1+|\Im z_1|+|\Im z_2|)^{-J_1}.
\end{equation}
To see why differentiation causes no arithmetic loss, each derivative
of $\rho/(y_1y_2)$ contributes a bounded multiple of $\rho$ on the
fixed compact support. For any fixed number of such derivatives, take
that many additional integrations by parts in $x$. The powers of
$\rho$ are then absorbed by its Fourier decay. All other factors are
bounded functions of $y_1,y_2$.

Mellin inversion first on right lines expresses the $\nu$ part of the
mode as
\[
 \frac1{\sqrt r(2\pi i)^2}
 \int\!\int \widetilde V_h(z_1,z_2)N_1^{z_1}N_2^{z_2}
 \mathcal G_{h,r}^{\nu}(1+it_1+z_1,1+it_2+z_2)
 \,dz_1\,dz_2.
\]
Move both lines to $\Re z_i=-1/2+b$. Lemma~\ref{analytic:euler}
and the entire continuation of the twist functions show that no pole
is crossed. Polynomial growth in fixed strips, together with a larger
fixed order in \eqref{analytic:mellin-kernel}, makes the horizontal
segments tend to zero. The scale factor is
$(N_1N_2)^{-1/2+b}$. Applying
\eqref{analytic:euler-bound} and \eqref{analytic:mellin-kernel}, and
discarding $r^{-1/2}\leq1$, proves
\eqref{analytic:nonzero-mode}. The final assertion follows directly
from the support of $\widehat W$ before Mellin inversion.
\end{proof}

Whenever unrestricted positive indices occur, write $n=2^a n'$ with
$n'$ odd. For $z_m(2)=c$ the factor taken outside an odd-index sum is
$\lambda_F(2^a)c^a2^{-a(1/2+it)}$. Its absolute values have the
uniformly convergent majorant
\begin{equation}\label{analytic:two-sum}
 \sum_{a\geq0}(a+1)2^{-a/2}<\infty.
\end{equation}
The new length is $N/2^a$; a nonempty localized sum has this length
bounded below by a constant depending only on $[a_0,a_1]$.

\subsection{Closing a logarithmic second-moment induction}

\begin{proof}[Proof of Propositions~\ref{analytic:moment} and~\ref{analytic:polynomial-moment}]
We give a simultaneous induction for all $F\in\mathcal F$.
The induction hypothesis at smaller dyadic heights is
\begin{equation}\label{analytic:induction-hypothesis}
 \frac1{M'}\sum_{\substack{d\ \mathrm{admissible}\\M'\leq|d|<2M'}}
 |L_{F^{(d)}}(\sigma+it)|^2
 \leq K\log^B(3M')(1+|t|)^A \qquad(M'<M).
\end{equation}
The constant $K$, common to the finite family, will be selected last.
Since the reduction in Lemma~\ref{setup:reduction} does not increase
height and has bounded multiplicity, this hypothesis implies
\begin{equation}\label{analytic:inductive-prefix}
 \sum_{\substack{0<|h_1|\leq H\\h_1\ \mathrm{squarefree}}}
 |L_{F^{(\nu h_1)}}(1/2+b+iv)|^2
 \leq C_{\mathcal F}KH\log^B(3M)(1+|v|)^A
 \quad(1\leq H\leq M/2).
\end{equation}
Indeed all admissible images lie in earlier dyadic blocks, whose
lengths have sum at most $2H$. If $H<1$, the sum is empty. In
particular the arguments below at $M=1$ require no induction input.

\smallskip
\noindent\emph{The polynomial average at inflated height.}
First we prove a polynomial estimate with a small coefficient in
front of $K$. Put $L=\log(3M)$, $T=1+|t|$, and let $N>0$.
Choose $P$ sufficiently large, with
\begin{equation}\label{analytic:inflation-condition}
 P\geq C_{\mathcal F,G}(1+N/M).
\end{equation}
Let $\mathcal P$ consist of the primes $p\in[P,2P]$ outside $S$.
For a sufficiently large fixed lower threshold on $P$, the elementary
dyadic prime estimate of Erd\H{o}s \cite[\S6, equation~(10)]{Erdos1932}
gives $|\mathcal P|\gg P/\log(2P)$.

Choose a fixed nonnegative Schwartz function $\Phi$ that is at least
one on $[-8,8]$ and has compactly supported Fourier transform. For
completeness, take a nonzero even real smooth bump $g$ in frequency
space with $\int g>0$. Its inverse Fourier transform $f$ is real,
Schwartz, and nonzero near zero. A sufficiently small dilation makes
$f(\alpha x)$ nonzero throughout $[-8,8]$; a fixed multiple of
$f(\alpha x)^2$ then has the required properties. Its Fourier
transform is a compactly supported convolution.

Put $U=MP^2$. For $c=\pm1$ use the formal character $z_m$ with
$z_m(2)=c$, and let $P_m^c(N',t;G)$ be the corresponding polynomial.
We claim, uniformly for $0<N'\leq N$,
\begin{align}
 \frac1U\sum_{m\in\mathbb Z}\Phi(m/U)|P_m^c(N',t;G)|^2
 &\leq C_G\log^{C_0}(3M+3N)
 \left(1+C_{A,B,\mathcal F,G}K\frac{N}{MP^2}L^BT^A\right),
 \label{analytic:inflated-moment}
\end{align}
where we may fix $C_0=512$, independently of $A,B$ and of the required
smoothness orders.

To prove the claim, apply Lemma~\ref{analytic:poisson} with $r=1$ to
the expansion of the square, first separating powers of two by
\eqref{analytic:two-sum}. Take
$b=1/(10\log(3M+3N))$. For the zero mode the complete-residue mean
vanishes unless the odd part of $n_1n_2$ is a square, and its absolute
value is at most one. Inserting $(n_1n_2)^{-b}$ at the bounded cost
$(C_GN)^{2b}\ll_G1$ majorizes the resulting sum by an Euler product.
At an odd prime this product has factor
\[
 \sum_{\substack{a,b'\geq0\\a+b'\ \mathrm{even}}}
       (a+1)(b'+1)p^{-(a+b')(1/2+b)}
 =1+O(100p^{-1-2b}).
\]
For example this follows by putting $z=p^{-1/2-b}\leq3^{-1/2}$ in
$\tfrac12((1-z)^{-4}+(1+z)^{-4})$ and bounding its terms of degree
at least two. The prime two costs a bounded factor by
\eqref{analytic:two-sum}. The zero mode is therefore $O_G(b^{-100})$.

For the nonzero modes, denote the two odd lengths by $N_1,N_2$.
The Fourier support restricts the frequencies to
\[
 0<|h|\leq H_0=C_\Phi N_1N_2/(MP^2).
\]
Increasing the constant in \eqref{analytic:inflation-condition}
ensures $H_0\leq M/2$ for every such pair of lengths. For fixed
$h_2$, Cauchy--Schwarz and \eqref{analytic:inductive-prefix} bound
the sum over $h_1$ occurring in \eqref{analytic:nonzero-mode} by
\[
 \ll_{\mathcal F}K\frac{H_0}{h_2^2}L^B
 (1+|t+v_1|)^{A/2}(1+|-t+v_2|)^{A/2}.
\]
This formula is only used when $H_0/h_2^2\geq1$; otherwise that sum
is empty. Choose $J_1>A+4$. Since
$1+|\pm t+v|\leq T(1+|v|)$, the integral in $v_1,v_2$ is
$O_A(T^A)$. Also
\begin{equation}\label{analytic:square-factor-sum}
 \sum_{h_2\geq1}\frac{\tau(h_2)^{C_{\rm ar}}}{h_2^2}<\infty.
\end{equation}
One proof of this convergence is its Euler product: its local factor
is $1+O_{C_{\rm ar}}(p^{-2})$, since
$\sum_{a\geq1}(a+1)^{C_{\rm ar}}p^{-2a}
\ll_{C_{\rm ar}}p^{-2}$ uniformly for $p\geq2$.
Consequently the nonzero contribution is bounded by
\[
 \ll_{A,B,\mathcal F,G}b^{-C_{\rm ar}}KL^BT^A
 (N_1N_2)^{-1/2+b}H_0
 \ll_{A,B,\mathcal F,G}
 \log^{C_0}(3M+3N)K\frac{N}{MP^2}L^BT^A.
\]
Here $(N_1N_2)^b\ll_G1$ and
$\sqrt{N_1N_2}\ll_GN$. The sums of the factors from two converge by
\eqref{analytic:two-sum}. This proves
\eqref{analytic:inflated-moment}. If the Fourier support contains no
nonzero frequency, the same proof uses only its zero-mode part.

\smallskip
\noindent\emph{Recovering the squarefree average.}
Now inflate an admissible parameter by setting $m=dp^2$. If
$c=\chi_d(2)$, multiplicativity gives the exact identity
\begin{equation}\label{analytic:inflation-identity}
 P_d(N,t;G)=\sum_{a\geq0}
 \lambda_F(p^a)\chi_d(p)^a p^{-a(1/2+it)}
 P_{dp^2}^c(N/p^a,t;G).
\end{equation}
Indeed $z_{dp^2}(n)$ deletes all positive powers of $p$ and agrees
with $\chi_d(n)$ on integers prime to $p$; reinstating their Euler
coefficients gives the original sum. If $p\mid d$, the terms with
$a\geq1$ vanish, so the identity also covers that case.

Equip pairs $(d,p)$, with $M\leq|d|<2M$, $\chi_d(2)=c$, and
$p\in\mathcal P$, with squared norm
$(M|\mathcal P|)^{-1}\sum_{d,p}|\cdot|^2$.
The map $(d,p)\mapsto dp^2$ is injective: in the prime factorization
of its absolute value the only exponent exceeding one occurs at
$p$, and is either two or three. Moreover $|dp^2|<8U$.
Thus the $a=0$ term has norm at most
\[
 \left(\frac{U}{M|\mathcal P|}\right)^{1/2}
 \sup_{0<N'\leq N}
 \left(\frac1U\sum_m\Phi(m/U)|P_m^c(N',t;G)|^2\right)^{1/2}
 \ll(P\log(2P))^{1/2}\mathcal B^{1/2},
\]
where $\mathcal B$ denotes the right side of
\eqref{analytic:inflated-moment}.
For $a\geq1$ fix $p$ before summing in $d$. The weaker bound
$\sum_d|P_{dp^2}^c(N/p^a,t;G)|^2\leq U\mathcal B$ gives pair norm
at most
$(a+1)P^{1-a/2}\mathcal B^{1/2}$. Minkowski's inequality and
\[
 \sum_{a\geq1}(a+1)P^{-a/2}\ll P^{-1/2}
\]
show that these terms together have norm
$O(P^{1/2}\mathcal B^{1/2})$. Summing the two choices of $c$ proves
\begin{align}
 \frac1M\sum_d|P_d(N,t;G)|^2
 &\ll_G\log^{C_0}(3M+3N)P\log(2P)
 \left(1+C_{A,B,\mathcal F,G}K\frac{N}{MP^2}L^BT^A\right).
 \label{analytic:deinflated-moment}
\end{align}

Here is an explicit noncircular choice of exponents. Set
\begin{equation}\label{analytic:exponent-choices}
 \begin{split}
 C_0&=512,\qquad C_1=C_0+32,\qquad C_2=C_0+C_1+4,\\
 R&=C_2+4,\qquad C_3=3C_2+10,\qquad
 A=B=2(C_3+20).
 \end{split}
\end{equation}
The letter $R$ in this proof denotes a fixed contour abscissa; it is
unrelated to a rank truncation threshold used later in the paper.
For any prescribed $\varepsilon>0$, choose
\begin{equation}\label{analytic:choice-P}
 P=P_0[L(1+N/M)T]^{C_1}.
\end{equation}
Once $A,B$ and the fixed weight family have been selected, a sufficiently
large $P_0$ ensures \eqref{analytic:inflation-condition} and gives
\begin{align}
 \frac1M\sum_d|P_d(N,t;G)|^2
 &\leq C_*[L(1+N/M)T]^{C_2}
 \nonumber\\
 &\quad+\varepsilon K L^{B-10}T^{A-4}(1+N/M)^{-4}.
 \label{analytic:small-recursion}
\end{align}
The constant $C_*$ can depend on $P_0$ and the chosen fixed data, but
not on $K,M,N,t$.
For explicit verification, put $q=1+N/M$ and $D_0=LqT$.
Then $\log(3M+3N)\ll Lq$ and
\[
 \frac{\log(2P)}P
 \leq\frac{\log(2P_0)+C_1}{P_0}D_0^{-C_1+1}.
\]
The recursive coefficient in \eqref{analytic:deinflated-moment} is
therefore at most a fixed multiple of
\[
 \frac{\log(2P_0)+C_1}{P_0}
 K L^{B+C_0-C_1+1}T^{A-C_1+1}q^{C_0-C_1+2}.
\]
Our value of $C_1$ leaves at least the losses $10,4,4$ displayed in
\eqref{analytic:small-recursion}. Its prefactor can be made smaller
than $\varepsilon$. The nonrecursive term is bounded by
$C_*D_0^{C_0+C_1+1}$ and hence by the stated $C_2$ power. In
particular $C_2$ was fixed before $A,B,\varepsilon,P_0,K$.

\smallskip
\noindent\emph{Reconstructing the central-line values.}
We next reconstruct $L$ on the central line. Choose a fixed \mbox{smooth}
dyadic partition $G_0$ with common compact support in $(0,\infty)$,
so that
$\sum_{N=1,2,4,\ldots}G_0(n/N)=1$ for every integer $n\geq1$.
Let $C_d=\sqrt{q_F}|d|/(2\pi)$ and define
\[
 I_d(t)=\frac1{2\pi i}\int_{(R)}
 L_{F^{(d)}}(1/2+it+w)C_d^w
 \frac{\Gamma(1+it+w)}{\Gamma(1+it)}\frac{e^{w^2}}w\,dw.
\]
Moving the contour to $\Re w=-R$ crosses only the residue
$L_{F^{(d)}}(1/2+it)$, because the completed function is entire.
Apply its functional equation on the new line and replace $w$ by
$-w$. One obtains the exact approximate functional equation
\begin{equation}\label{analytic:afe-identity}
 L_{F^{(d)}}(1/2+it)=I_d(t)+\omega_d(t)I_d(-t),\qquad
 \omega_d(t)=\epsilon_d C_d^{-2it}
                       \frac{\Gamma(1-it)}{\Gamma(1+it)},
 \quad |\omega_d(t)|=1.
\end{equation}
The Gaussian makes the contour displacement legitimate. Initially,
on $\Re w=R$, the Dirichlet series is absolutely convergent, so we
may insert the dyadic partition.

For $0\leq c\leq R$, Stirling's formula uniformly gives
\begin{equation}\label{analytic:gamma-quotient}
 \left|\frac{\Gamma(1+c+i(t+v))}{\Gamma(1+it)}\right|
 \ll_R T^c(1+|v|)^{R+1/2}e^{\pi|v|/2}.
\end{equation}
Indeed the quotient of the exponential factors is at most
$e^{\pi|v|/2}$, and
\[
 \frac{(1+|t+v|)^{c+1/2}}{(1+|t|)^{1/2}}
 \leq T^c(1+|v|)^{c+1/2}.
\]
The same estimates with $1+|t+v|$ also cover bounded ordinates, including
$v$ near $-t$. Thus the bound is uniform as $c$ tends to zero.
For a fixed dyadic piece the polynomial is finite, so its contour can
be moved to any $0<c\leq R$. Multiplying
\eqref{analytic:gamma-quotient} by $|e^{(c+iv)^2}|$ absorbs its
$v$-factors into $O_R(e^{-v^2/2})$ and yields
\begin{equation}\label{analytic:afe-pieces}
 |L_{F^{(d)}}(1/2+it)|
 \ll_{R,\mathcal F}\sum_{\nu=\pm1}\sum_N
 \int_{\mathbb R}\frac{e^{-v^2/2}}{|c_N+iv|}
      \left(\frac{MT}N\right)^{c_N}
 |P_d(N,\nu t+v;G_0(y)y^{-c_N})|\,dv.
\end{equation}
There is no use of the desired moment to justify this representation:
on the initial fixed line $R$, the elementary bound
$|P_d(N,t;G)|\ll_G\sqrt N\log(2+N)$ makes the sum of the long
pieces converge, since $R>1/2$. Finitely many shorter pieces can
then be moved separately. The moment estimates below also show
convergence after taking the family norm.

Use
\[
 c_N=\frac1{\log(3MT)}\quad(N\leq MT^2),\qquad
 c_N=R\quad(N>MT^2).
\]
Apply \eqref{analytic:small-recursion} uniformly to the fixed smooth
family $G_0(y)y^{-c}$, $0\leq c\leq R$, and use Minkowski in
\eqref{analytic:afe-pieces}. We spell out the two contributions.
For the nonrecursive square-root term and $N\leq MT^2$,
\[
 [L(1+N/M)(1+|\nu t+v|)]^{C_2/2}
 \ll L^{C_2/2}T^{3C_2/2}(1+|v|)^{C_2/2}.
\]
There are $O(\log(3MT^2))$ such dyadic pieces,
$(MT/N)^{c_N}\leq e$ for $N\geq1$, and
$1/|c_N+iv|\leq\log(3MT)$.
Since both logarithms are $O(LT)$, their squared total cost gives
at most $C_{**}(LT)^{3C_2+4}$ in the mean square.
For the long pieces put $q'=N/M>T^2$. Then
\[
 \sum_{q'>T^2}\left(\frac T{q'}\right)^R(q')^{C_2/2}
 \ll_R T^{-R+C_2},
\]
where $q'$ runs over a dyadic progression. Including the outside
factor $L^{C_2/2}T^{C_2/2}$ again fits within
$C_{**}(LT)^{C_3}$ by \eqref{analytic:exponent-choices}.

For the recursive square-root term, the short dyadic sum satisfies
\[
 \sum_{N\leq MT^2}(1+N/M)^{-2}\ll\log(3M)=L.
\]
The Gaussian absorbs $(1+|v|)^{(A-4)/2}$, so the resulting norm is
at most
\[
 C_{A,R,\mathcal F}\sqrt{\varepsilon K}\,
 L^{B/2-5}T^{A/2-2}\,L\log(3MT)
 \ll C_{A,R,\mathcal F}\sqrt{\varepsilon K}\,
 L^{B/2}T^{A/2}.
\]
For the long pieces we have the stronger estimate
\begin{equation}\label{analytic:long-recursive-sum}
 \sum_{q'>T^2}\left(\frac T{q'}\right)^R(1+q')^{-2}
 \ll_R T^{-R-4},
\end{equation}
which gives the same bound. Squaring the total norm now proves
\begin{equation}\label{analytic:central-recursion}
 \frac1M\sum_d|L_{F^{(d)}}(1/2+it)|^2
 \leq C_{**}(LT)^{C_3}
       +C_{\dagger}\varepsilon K L^BT^A.
\end{equation}
Here $C_{\dagger}$ depends on $A,R,\mathcal F$ but not on $P_0,K,M,t$;
$C_{**}$ may depend on $P_0$. This independence is essential when
choosing the constants.

\smallskip
\noindent\emph{Closing the induction in the strip.}
It remains to extend the estimate to the strip without a factor for
the number of twists. Form the vector
\[
 \mathbf V_F(s)=\big(M^{-1/2}L_{F^{(d)}}(s)\big)_d
\]
in the finite-dimensional Euclidean space indexed by the current
admissible block. On $\Re s=2$ its squared norm is bounded by an
absolute constant: each $L$-value has modulus at most
$\sum_n\tau(n)n^{-2}=\zeta(2)^2$, and the number of parameters is
$O(M)$. On $\Re s=1/2$, \eqref{analytic:central-recursion} gives the
stated vector norm bound.

For each constant unit vector $u$ apply the scalar strip principle to
\[
 f_u(s)=\frac{\langle\mathbf V_F(s),u\rangle}{(1+s)^{A/2}}.
\]
Here $A$ is even by \eqref{analytic:exponent-choices}. On both
boundaries $|1+s|\asymp1+|\Im s|$, with constants independent of the
number of twists. The scalar boundary bound is therefore at most
\[
 C_A\big(C_{**}L^{C_3}+C_{\dagger}\varepsilon K L^B+1\big)^{1/2}.
\]
For clarity, finite-order growth suffices for the strip principle:
multiply $f_u(s)$ by
$\exp\{-\rho\cos(s-5/4)\}$, with $\rho>0$. Its modulus is at most
one on the strip boundaries and it decays faster than every
finite-order growth bound on horizontal sides, since
$\cos(\Re s-5/4)>0$ for $1/2\leq\Re s\leq2$. Apply the maximum
principle on rectangles and then let their heights tend to infinity
and $\rho$ decrease to zero. Taking the supremum over unit vectors
at the desired point gives
\begin{equation}\label{analytic:strip-recursion}
 \|\mathbf V_F(\sigma+it)\|_2^2
 \leq C_{\rm st}
       \big(C_{**}L^{C_3}+C_{\dagger}\varepsilon K L^B+1\big)T^A,
 \qquad 1/2\leq\sigma\leq2.
\end{equation}
This argument has no dimension-dependent loss.

The choice order is now explicit. Fix the family, the elementary
arithmetic exponent, $C_0,C_1,C_2,R,C_3,A,B$ as above, and then all
needed fixed kernel orders. Next choose
$0<\varepsilon<(4C_{\rm st}C_{\dagger})^{-1}$, choose $P_0$ to obtain
\eqref{analytic:small-recursion} for the entire weight family, and
finally choose $K$ so large that
$C_{\rm st}(C_{**}L^{C_3}+1)\leq K L^B/2$ for every $L\geq\log3$.
This is possible because $B>C_3$ and $C_{**}$ is independent of $K$.
Equation~\eqref{analytic:strip-recursion} proves
\eqref{analytic:induction-hypothesis} at height $M$. At $M=1$, all
dual sums used above are empty, so this also establishes the initial
step. The simultaneous induction is complete.

The induction just completed uses the fixed weight family
$G_0(y)y^{-c}$, $0\leq c\leq R$. For any other fixed smooth compactly
supported $G$, repeat the polynomial argument through
\eqref{analytic:small-recursion}, using the established $L$-function
bound at smaller heights. The constants $P_0$ and $C_*$ may now depend
on $G$, while $K,A,B$ remain fixed. Since $A,B>C_2$, this proves
\eqref{analytic:polynomial-bound}, with the asserted uniformity for
fixed uniformly smooth weight families.
Finally use Lemma~\ref{setup:reduction}, bounded multiplicity, and a
sum over admissible dyadic blocks of total length $O(H)$ to obtain
\eqref{analytic:prefix-moment}. No analytic-density input has entered
the argument.
\end{proof}

\subsection{Short polynomials and auxiliary multipliers}

The second moments are now available at every height. We use them to
bound the complete sum of nonzero Poisson frequencies, retaining the
zero frequency as an exact independent-model average. The resulting
comparison permits a short polynomial multiplier, as required for the
mollifiers below.

Fix a nonnegative smooth function $\Psi$ compactly supported in
$\{x:0<|x|<\infty\}$, with $\Psi(x)\geq1$ when $1\leq|x|\leq2$.
Fix $c\in\{1,-1\}$. The formal character $z_m$ and the independent
variables $Z$ always have value $c$ at two. At odd primes their model
law is
\[
 \mathbb P(Z(p)=0)=\frac1p,\qquad
 \mathbb P(Z(p)=1)=\mathbb P(Z(p)=-1)=\frac{1-1/p}{2}.
\]
They are extended completely multiplicatively. This is a model for
complete integer residue averages.

\begin{lemma}[Summation of all dual frequencies]\label{analytic:frequency-sum}
For fixed $B\geq0$, $J_1>1$, and $Q>0$,
\begin{equation}\label{analytic:dyadic-frequency-bound}
 \sum_{T=1,2,4,\ldots}T\log^B(3T)(1+T/Q)^{-J_1}
 \ll_{B,J_1}
 \begin{cases}
 Q^{J_1},&0<Q<1,\\
 Q\log^B(3Q),&Q\geq1.
 \end{cases}
\end{equation}
In particular the first bound is at most a constant times $Q$.
\end{lemma}

\begin{proof}
If $Q<1$, bound $(1+T/Q)^{-J_1}$ by $(Q/T)^{J_1}$.
The remaining series is
$\sum_{j\geq0}2^{-(J_1-1)j}\log^B(3\cdot2^j)<\infty$.
For $Q\geq1$, the part $T\leq Q$ is bounded by
$\log^B(3Q)\sum_{T\leq Q}T\ll Q\log^B(3Q)$.
For the remaining part choose the first dyadic $T_0>Q$, so
$Q<T_0\leq2Q$, and write $T=2^jT_0$. The summand is at most
\[
 C_B Q\,2^{-(J_1-1)j}
       \big(\log^B(3Q)+(j+1)^B\big).
\]
Both resulting geometric series converge. This estimates the infinite
frequency tail before any logarithm is replaced by $\log X$.
\end{proof}

\begin{proposition}[Comparison with the independent model]
\label{analytic:comparison}
Let $\ell=\log X$, where $X$ is sufficiently large, and let
\[
 c_1\ell^{-1/10}\leq\Delta\leq\frac12
\]
for a fixed $c_1>0$. There are fixed constants $\eta,c_2>0$ for which
the following holds. Let
\[
 S_V(x)=\sum_{n\geq1}\frac{\lambda_F(n)V(n)}{\sqrt n}\,
                         w_x(\ell^{-1}\log n),\qquad
 H_V=\sum_{r'\leq R_0}h(r')V(r'),
\]
where $V=z_m$ or $Z$. Suppose that $S_V$ has length at most
$X^{1-\Delta}$, that
\begin{equation}\label{analytic:multiplier-conditions}
 1\leq R_0\leq X^{\eta\Delta},\qquad
 \sum_{r'\leq R_0}|h(r')|\leq C_H R_0^2,
\end{equation}
and that $w_x$ and a function $A(x)$ have bounds by a fixed power of
$\ell$ for all derivatives up to a sufficiently large fixed order.
For $w_x$ these bounds are required after smooth dyadic localization
in $n$, on the scaled variables $n/N$ and $x$ in a fixed neighborhood
of $\operatorname{supp}\Psi$. They may in particular follow from
bounds in $u=\ell^{-1}\log n\geq0$ and $x$, with a smooth extension
at $n=1$. Then
\begin{align}
 &\frac1X\sum_{m\in\mathbb Z}\Psi(m/X)
       |H_{z_m}S_{z_m}(m/X)-A(m/X)|^2
 \nonumber\\
 &\hspace{8mm}=
 \int_{\mathbb R}\Psi(x)\,
       \mathbb E|H_ZS_Z(x)-A(x)|^2\,dx
       +O\big(\exp(-c_2\Delta\ell)\big).
 \label{analytic:model-comparison}
\end{align}
The constants and the sufficiently large $X$ threshold are uniform
when the stated support, coefficient, and derivative bounds are
uniform. They do not depend on $\Delta$ or on an index parametrizing
such weights. One may take
\begin{equation}\label{analytic:comparison-constants}
 C'=C_{\rm ar}+1=257,\qquad
 \eta=\frac1{4(4+2C')},\qquad c_2=\frac12.
\end{equation}
\end{proposition}

\begin{proof}
Put $L=X^{1-\Delta}$. If the length hypothesis specifies only that
the integer coefficients vanish for $n>L$, multiply the interpolating
weight by $\rho_+(n/L)$, where $\rho_+$ is a fixed smooth function equal
to one on $(-\infty,1]$ and zero on $[2,\infty)$. This preserves every
integer coefficient and gives real support in $n\leq2L$. On a dyadic
scale meeting the transition, $N/L$ is bounded, so derivatives of
$\rho_+(Ny/L)$ add only fixed constants to the assumed scaled seminorm
bounds. Thus all nonempty localized lengths satisfy $N_i\ll L$.

Expand the square on the left of \eqref{analytic:model-comparison}.
A term involving two multiplier indices $r_1',r_2'$ has character
factor $z_m(r_1'r_2'n_1n_2)$. Its odd multiplier part is an integer
$r\leq R_0^2$. Powers of two contribute only fixed signs. Split
powers of two out of $n_1,n_2$, and insert smooth dyadic partitions
in their odd parts. Each resulting term has exactly the form
\eqref{analytic:poisson-average} with $U=X$, $t_1=t_2=0$, and a
weight $W$ whose required seminorm is $O(\ell^{D_1})$ for a fixed
$D_1$. The independent factors coming from the powers of two satisfy
\eqref{analytic:two-sum}. Cross terms with $A(x)$ are covered by
using an index supported just at one: choose a fixed smooth function
supported in $(1/2,3/2)$ and equal to one at one. The identity
$\lambda_F(1)=1$ then puts them in the same form. The term $|A|^2$
is covered by two such indices. Its derivatives obey the same
fixed polynomial bounds.

The zero frequency is exactly the model expression. In fact, for
an odd prime and $v\geq1$,
\[
 \frac1{p^v}\sum_{a\bmod p^v}\left(\frac a{p^v}\right)
 =\begin{cases}0&v\text{ odd},\\1-p^{-1}&v\text{ even},\end{cases}
 =\mathbb E Z(p)^v.
\]
For $v=0$ both sides are one. Chinese remaindering gives independence
at distinct odd primes, and the value at two is already fixed.
The integral of the zero-frequency smooth amplitude is thus the
corresponding term on the right of
\eqref{analytic:model-comparison}.

It remains to estimate the nonzero frequencies. Fix an odd multiplier
$r$ and localized odd lengths $N_1,N_2$, and set
\[
 b=\frac1{10\ell},\qquad Q=\frac{rN_1N_2}{X}.
\]
Only nonempty lengths need be considered, so $N_i$ is bounded below
by a fixed positive constant. Also $N_i\ll X^{1-\Delta}$, whence
$(N_1N_2)^b\ll1$. Proposition~\ref{analytic:moment} is available
at every height. For a dyadic frequency range $T\leq|h|<2T$,
write $h=h_1h_2^2$. For fixed $h_2$, the parameter $h_1$ lies in
$0<|h_1|<2T/h_2^2$. Cauchy--Schwarz and
\eqref{analytic:prefix-moment} imply
\begin{align*}
 &\sum_{\substack{T\leq|h|<2T\\h=h_1h_2^2}}
 \tau(h_2)^{C_{\rm ar}}
 \prod_{i=1}^2
 |L_{F^{(\nu h_1)}}(1/2+b+iv_i)|\\
 &\quad\ll
 T\log^B(3T)(1+|v_1|)^{A/2}(1+|v_2|)^{A/2}
 \sum_{h_2\leq\sqrt{2T}}\frac{\tau(h_2)^{C_{\rm ar}}}{h_2^2}\\
 &\quad\ll T\log^B(3T)
               (1+|v_1|)^{A/2}(1+|v_2|)^{A/2}.
\end{align*}
This inequality applies separately to each $\nu$; it includes both
signs of $h$. Choose $J_1>A+4$ in
\eqref{analytic:nonzero-mode}, and then fix the corresponding
smoothness order. Its vertical integral is convergent. The remaining
sum is bounded, by Lemma~\ref{analytic:frequency-sum}, by
\[
 \ll \ell^{D_2}r^{C_{\rm ar}}(N_1N_2)^{-1/2}
 \begin{cases}Q^{J_1}&Q<1,\\Q\log^B(3Q)&Q\geq1.
 \end{cases}
\]
When $Q\geq1$ we have $Q\ll X$ using $r\leq R_0^2$,
\eqref{analytic:multiplier-conditions}, and $\eta<1$; hence
$\log(3Q)\ll\ell$. When $Q<1$ use $Q^{J_1}\leq Q$.
In both regimes the error for this localized pair is therefore
\begin{equation}\label{analytic:localized-comparison-error}
 \ll\ell^{D_3}r^{C'}\frac{\sqrt{N_1N_2}}X,
 \qquad C'=C_{\rm ar}+1.
\end{equation}
The increase from $C_{\rm ar}$ to $C'$ is only the factor $r$
in $Q$, not a smoothness loss. Notice that this argument controls
arbitrarily large dual frequencies as well as the case $Q<1$.

There are $O(\ell)$ possible nonempty dyadic pieces for each main
index. Their cost is a fixed logarithmic power. Summing the
coefficients from powers of two by \eqref{analytic:two-sum} costs
only a constant. The two multiplier coefficient sums and the
bound $r\leq R_0^2$ cost at most
\[
 (1+\textstyle\sum|h(r')|)^2R_0^{2C'}\ll_{C_H}R_0^{4+2C'}.
\]
Since $\sqrt{N_1N_2}/X\ll X^{-\Delta}$, including the terms with an
index at one, the full error is
\begin{equation}\label{analytic:total-comparison-error}
 \ll\ell^{D_4}X^{-\Delta}R_0^{4+2C'}
 \leq\ell^{D_4}\exp(-3\Delta\ell/4).
\end{equation}
Our choices in \eqref{analytic:comparison-constants} give the last
inequality. The power $D_4$ is fixed before choosing $X$.
Because $\Delta\ell\geq c_1\ell^{9/10}$, for one sufficiently large
$X$ the factor $\ell^{D_4}$ is at most
$\exp(\Delta\ell/4)$. This proves
\eqref{analytic:model-comparison} with $c_2=1/2$.
\end{proof}

The endpoint at $n=1$ in this proposition needs no extension to a fixed
negative interval in $u$.
Extend a dyadically localized weight only to, say, $n\geq1/2$.
In the coordinate $\log n=\ell u$ this is a fixed interval, and
smooth cutoffs there introduce at most fixed powers of $\ell$ in
$u$ derivatives. On a dyadic scale, repeated differentiation of
$w_x(\log(Ny)/\ell)$ instead introduces powers of $1/\ell$ and
bounded powers of $1/y$. Thus the required finite seminorm bounds
are precisely those in the proposition, even at the first piece.
